\scanpage{160}
En particulier, si $X=Y$, $F=G$, prenant $\eta^0=\deleted{\unsure{\alpha}}\delta_{X,F}$ on trouve
\[
\boxed{\operatorname{Tr}(\xi^0\cup\delta_{X,F}^{0})=\sum\unsure{(-1)^i}\operatorname{Tr}[\xi^0]^i}
\]
C'est la \uline{formule de Lefschetz-Verdier}.

\margin{\textcircled{e}\quad 2}
\deleted{\unsure{\S}}
Considérons un morphisme
\[
u:T\longrightarrow X\times_S Y
\]
\unsure{défini} donc par deux morphismes
\[
u_1:T\longrightarrow X,\quad u_2:T\longrightarrow Y.
\]
On a alors un morphisme évident
\[
\mathbb{R}u_*\bigl(\mathbb{R}\mathcal{H}om(u_1^*(F_\bullet),u_2^!(G_\bullet))\bigr)
\longrightarrow\mathbb{R}\mathcal{H}om(\mathrm{pr}_1^*(F),\mathrm{pr}_2^!(G_\bullet))
\]
Car $u_i=\mathrm{pr}_i u$, donc le premier est un
\diagram{\begin{tikzcd}[ampersand replacement=\&,column sep=large,row sep=large]
\mathbb{R}u_*\bigl(\mathbb{R}\mathcal{H}om(\deleted{\unsure{\mathrm{pr}}}u^*\mathrm{pr}_1^*(F),u^!(\mathrm{pr}_2^!(G_\bullet)))\bigr)\arrow[r,"\sim"]\arrow[d,phantom,"\simeq" description]
\&\mathbb{R}\mathcal{H}om(\mathrm{pr}_1^*(F),u_*u^!(\mathrm{pr}_2^!(G)))\arrow[d,"\varepsilon"]\\
\mathbb{R}\mathcal{H}om(u_*u^*(\mathrm{pr}_1^*(F)),\mathrm{pr}_2^!(G_\bullet))\arrow[r]
\&\mathbb{R}\mathcal{H}om(\mathrm{pr}_1^*(F),\mathrm{pr}_2^!(G))
\end{tikzcd}}

\deleted{Donc} tout élément de \deleted{$\unsure{\operatorname{Corr}^{0}(X,F,)}$}
\[
\alpha^0\in\mathbb{R}\operatorname{Hom}^{0}(u_1^*(F),u_2^!(G))
=\mathbb{R}^{0}\Gamma_T\bigl(\mathbb{R}\mathcal{H}om(u_1^*(F),u_2^!(G_\bullet))\bigr)
\]
définit un \unsure{élément}
\[
\overline{\alpha^0}\in\operatorname{Corr}(X,F;Y,G)
\]
De même, un élément $\beta^0$ de $\mathbb{R}\operatorname{Hom}^{0}(u_1^*(F),u_2^!(G))$ définit un \unsure{élmt} $\overline{\beta}\in\operatorname{Corr}(Y,G;X,F)$.
\margin{159}
